\documentclass{leveldoc}
\DocUses{foundations,slingshot,rf_kick}
\DocFingerprint{e1c3eb0ab54cb949}

\begin{document}

\section{What the level asks}\label{sec:asks}

The game's text:
\leveldescription

Two identical ions ($q = \val{q}$, $m = \val{m}$, $T_0 = \val{Tzero}$) are launched from
$(\val{xZero}, \val{yZero})$ along $+x$ in a world with $c = \val{c}$: the first at $t = 0$, the second
at $t_0 = \val{tLaunch}$, which is $\val{halfRatio}$ of the generator's period $2\pi/\omega$
($\omega = \val{omega}$). The first must enter the detector $\val{xLoOne} \le x \le \val{xHiOne}$,
$\val{yLoOne} \le y \le \val{yHiOne}$ above the beam line, the second
$\val{xLoTwo} \le x \le \val{xHiTwo}$, $\val{yLoTwo} \le y \le \val{yHiTwo}$ below it. The player
places up to two antennas, as in \levelref{rf_kick} (amplitudes from the list, either sign, four
orientations), and one charge, at nodes with $\val{xRegLo} \le x \le \val{xRegHi}$,
$\val{yRegLo} \le y \le \val{yRegHi}$. A flight ends when the ion enters its detector, touches an
antenna's body (radius $\val{Ra}$), leaves the world, or reaches $t = \val{tMax}$
(\xsecref{foundations}{sec:board}).

The reference places one antenna, $p_0 = \val{pA}$ at $(\val{xA}, \val{yA})$, oriented at
$\val{angleA}^\circ$. The document shows that the late ion is exactly the early one with its charge
reversed, so the antenna kicks the two oppositely to first order, marches both flights, and
measures the second-order difference.

\section{Setup}\label{sec:setup}

The level file, field by field (the appendix reproduces it):
\setuptable

By \xeqref{foundations}{eq:launch} each ion starts with $\gamma_0 = \val{gammaZero}$, momentum
$\val{pZero}$ and speed $v_0 = \val{vZero}$. The drive's period is $\val{period}$.

\begin{figure}[htbp]
\centering
\begin{tikzpicture}
\begin{axis}[axis equal image, width=0.92\linewidth, xmin=-1.6, xmax=31.6, ymin=-1.6, ymax=21.6,
  xtick={0,5,...,30}, ytick={0,5,...,20}, tick label style={font=\small},
  xlabel={$x$}, ylabel={$y$}, label style={font=\small}]
  \draw[black!8] (0,0) grid (30,20);
  \draw[black!45, dashed] (-1,-1) rectangle (\rawval{xWorld},\rawval{yWorld});
  \draw[black!30, dotted] (\rawval{xRegLo},\rawval{yRegLo}) rectangle (\rawval{xRegHi},\rawval{yRegHi});
  \fill[green!50!black, fill opacity=0.18] (\rawval{xLoOne},\rawval{yLoOne}) rectangle (\rawval{xHiOne},\rawval{yHiOne});
  \fill[green!50!black, fill opacity=0.18] (\rawval{xLoTwo},\rawval{yLoTwo}) rectangle (\rawval{xHiTwo},\rawval{yHiTwo});
  \addplot[thick, blue!70!black] table {pathone.dat};
  \addplot[thick, red!70!black] table {pathtwo.dat};
  \fill[red!75!black] (\rawval{xA},\rawval{yA}) circle[radius=\rawval{Ra}];
  \fill (\rawval{xZero},\rawval{yZero}) circle[radius=0.15];
\end{axis}
\end{tikzpicture}
\caption{The reference antenna (red disc), the placement region (dotted), the detectors and the
two flights: launched at $t = 0$ (blue) and half a period later (red).}
\label{fig:board}
\end{figure}

\section{Principles and model}\label{sec:model}

\begin{principles}
  \item The antenna's exact retarded fields, \xeqref{rf_kick}{eq:fields}, derived in
    \levelref{rf_kick} from Maxwell's equations.
  \item The relativistic equations of motion \xeqref{foundations}{eq:eom} with the launch
    \xeqref{foundations}{eq:launch}; Liénard's power \xeqref{foundations}{eq:larmor}.
  \item The kick to first order in the field, \xeqref{rf_kick}{eq:first-order}.
  \item The board's rules; the classical Runge--Kutta method, cubic Hermite interpolation and the
    error estimate from the step $2h$ (\xsecref{foundations}{sec:board},
    \xsecref{foundations}{sec:numerics}), as in \xsecref{slingshot}{sec:method}.
\end{principles}

An ion launched at $t_0$ moves in the fields of the lab time $t_0 + t$, $t$ counted from its own
launch (\physref{2.4}):
\begin{equation}\label{eq:offset}
  \dot{\vect{p}} = q\big(\vect{E}(\vect{x}, t_0 + t) + \vect{v}\times\vect{B}(\vect{x}, t_0 + t)\big),
  \qquad \vect{x}(0) = (\val{xZero}, \val{yZero}), \quad \vect{p}(0) = (p_\mathrm{i}, 0).
\end{equation}

\begin{assumption}{The ions' radiation, and the antenna as a fixed source}\label{as:rf}
As in \levelref{rf_kick}, whose assumptions \xref{rf_kick}{as:radiation} and
\xref{rf_kick}{as:driven} these are: the ions' radiation is neglected, $\val{radOne}$ and $\val{radTwo}$ here,
and the antenna oscillates at a fixed amplitude. The two ions fly separately and do not act on
each other. These are the game's model.
\end{assumption}

\section{Derivation}\label{sec:derivation}

\subsection{Half a period late is the opposite charge}\label{sec:conjugate}

The dipole $\vect{p}(t) = \vect{p}_0\cos\omega t$ satisfies $\vect{p}(t + \pi/\omega) = -\vect{p}(t)$,
and so do its derivatives; the fields \xeqref{rf_kick}{eq:fields} are linear in them, at the
retarded time $t - r/c$, so $\vect{E}(\vect{x}, t + \pi/\omega) = -\vect{E}(\vect{x}, t)$ and the same for
$\vect{B}$. With $t_0 = \pi/\omega$, \eqref{eq:offset} for the second ion reads
\begin{equation}\label{eq:conjugate}
  \dot{\vect{p}} = (-q)\big(\vect{E}(\vect{x}, t) + \vect{v}\times\vect{B}(\vect{x}, t)\big):
\end{equation}
the second ion moves exactly as the first would with its charge reversed. Marching both confirms it
to $\val{conjDiff}$ in the end state. The antenna cannot tell the two apart by anything but this
sign, which is what an RF separator uses: no static field reverses with the launch time.

\subsection{The kicks}\label{sec:kicks}

To first order in the field \xeqref{rf_kick}{eq:first-order} is linear in $q$, so by
\eqref{eq:conjugate} the two ions are lifted by equal and opposite amounts: integrated along the
undeflected line to $t^* = \val{tStar}$, when it reaches the detectors' face, $\pm\val{liftFirst}$.
The full flights differ from this at second order, where the field is sampled along the deflected
paths, and that part is even in $q$: it moves both ions the same way, so one is lifted by more
than the other is lowered, by twice that part.

\subsection{The march}\label{sec:march}

Each flight \eqref{eq:offset} is marched by \xeqref{foundations}{eq:rk4} at $h = \val{H}$ with
Liénard's energy as one more state; the entry is located on the cubic Hermite interpolant by
bisection and the error estimated against the step $2h$, as in \xsecref{rf_kick}{sec:flight}.

\section{The reference solution by hand}\label{sec:hand}

\begin{procedure}
  \item\label{st:one} \emph{The first ion.} The march of \S\ref{sec:march} with $t_0 = 0$: it enters
    its detector through the \val{faceOne} face at $(\val{xOne}, \val{yOne})$ at $t = \val{tOne}$, with
    $\vect{p} = (\val{pxOne}, \val{pyOne})$ (step estimate $\val{errOne}$).
  \item\label{st:two} \emph{The second ion.} With $t_0 = \val{tLaunch}$, or equally with its charge
    reversed and $t_0 = 0$ \eqref{eq:conjugate}: through the \val{faceTwo} face at
    $(\val{xTwo}, \val{yTwo})$ at $t = \val{tTwo}$ after its launch, with
    $\vect{p} = (\val{pxTwo}, \val{pyTwo})$ (step estimate $\val{errTwo}$).
  \item\label{st:asym} \emph{The asymmetry.} The first ion is lifted by $\val{liftUp}$, the second
    lowered by $\val{liftDown}$, against $\val{liftFirst}$ each to first order: the second-order part
    is $\val{asym}$, $\val{asymRel}$ of the lift.
\end{procedure}

\begin{result}[The reference flights]
Each ion enters its own detector: the first above the beam line, the second below. The level is
solved.
\end{result}

\section{Comparison with the game}\label{sec:comparison}

\comparisontable

All trajectory differences are at most $\val{largest}$, the game's verified accuracy (Foundations
§6). The tolerances are ten times the larger of the two computations' estimates, or the game's own
preview--verification gap where that is larger; the radiated energies agree to a few $10^{-9}$
relative. On the empty board both ions fly straight to the world's right edge at
$t = \val{bareT}$.

\section{Other solutions}\label{sec:others}

Any setup that solves \levelref{rf_kick}'s problem for one ion, lifting it into the upper
detector, sends the other ion down by about as much, by \eqref{eq:conjugate}; the detectors are
placed symmetrically about the beam line and three cells tall, so the second-order asymmetry
($\val{asymRel}$ of the lift here) must stay within their margins. The level's charge cannot help
by itself: its static field kicks both ions alike. A map of all solutions would take the game's
verdicts, which this document does not reproduce.

\section{What the game computes}\label{sec:engine}

The game evaluates the antenna's fields at the lab time $t_0 + t$ for a shot launched at $t_0$
(\physref{2.4}, exact) and integrates the motion with its adaptive eighth-order Runge--Kutta method
(\physref{3}, \physref{5}, \physref{6}, \physref{7}), each shot verified on its own. The problem is
the same as here, and the differences in the comparison are the two integrations' errors.

\section{Further reading}\label{sec:reading}

The antenna's fields and one ion's kick: \levelref{rf_kick}. The oscillating dipole:
\Jackson{§9.2}.

\levelfileappendix

\end{document}
